\section{History}

\begin{wrapfigure}[8]{r}{0.20\textwidth}
\centering
\includegraphics[width=.20\textwidth]{drawings/NeXT_logo.pdf}
\end{wrapfigure}
\par
NeXT's history starts (and amusingly, also ends) at Apple. In May of 1985, the Macintosh mediocre sales drew a bleak future for the company. Steve Jobs, co-founder and then General Manager of the Mac department, wanted to lower the price and increase marketing in order to boost the Mac. John Sculley then CEO, wanted to abandon the Mac and refocus the company resources on the Apple II, the only profitable product Apple had marketed until then.\\
\par
 A vote was called and the board of directors sided with Sculley. Steve Jobs found himself stripped of all responsibilities. A few months later, on September 13, 1985, he resigned and went on to work on his next project.\\
\par
NeXT, Inc was incorporated in February 1986 with \$7 million of Jobs's own money. Many member of the Mac division left Apple to join the newly formed company, among them  Joanna Hoffman, Bud Tribble (head of software division), George Crow, Rich Page, Susan Barnes, Susan Kare, and Dan'l  Lewin.\\ 
\par
With NeXT, Jobs went back to a project he had contemplated while still at Apple in August 1985. While touring universities to boost Mac sales, he had met Paul Berg, a Nobel Laureate in chemistry frustrated by the expense (\$100,000) of teaching students about recombinant DNA from textbooks instead of in wet laboratories. It seemed there was a market for powerful workstation targed at universities and students\footnote{The Second Coming of Steve Jobs.}. NeXT set itself to build a workstation powerful enough to run wet lab simulations and cheap enough for college students to use in their dormitory rooms.\\
\par
\fq{I want some kid at Stanford to be able to cure cancer in his dorm room.}{Steve Jobs, 1987}
\par
Steve Jobs spared no expenses. For \$100,000, Paul Rand was commissioned with a logo. He also decided to follow Alan Kay precept to make both the hardware and the software that would power his machine.\\
\par
\fq{People who are really serious about software should make their own hardware.}{Alan Kay, 1980}
\par
Using his experience on the Macintosh, he defined the three pillars NeXT Computer: GUI, Networking and Object-Oriented programming.\\
\par
\fq{I went to Xerox PARC. And they were very kind. They showed me what they are working on. And they showed me really three things. But I was so blinded by the first one that I didn't even really see the other two. One of the things they showed me was object oriented programming -- they showed me that but I didn't even see that. The other one they showed me was a networked computer system... they had over a hundred Alto computers all networked using email etc., etc., I didn't even see that. I was so blinded by the first thing they showed me, which was the graphical user interface. I thought it was the best thing I'd ever seen in my life. }{Steve Jobs, 1995}

\par
\section{The NeXT Computer}
The first machine shipped in 1989 after three years of hard work. Not meeting the initial release target was not a problem according to Jobs who famously replied to a journalist inquiring about the delay: "Late? This computer is five years ahead of its time!".\\
\par
Based on a Motorola 68030 25 Mhz and featuring powerful co-processors such a Digital Signal Processor and a Floating Point Unit, the hardware delivered. On top of a powerful architecture, the machines coming out of Red Woods, CA factory also happened to be gorgeous. The one foot perfect cube made of painted magnesium was elegant.\\
\par
\cfullimage{NEXT_Cube}{The Next Computer}


The monitor was a piece of art itself. The MegaDisplay allowed a resolution of 1120 x 832 pixels with a density of 92 DPI (just shy of a true Megapixel at 931,840 total pixels). At launch the supply chain was so tight that when ordering a NeXTComputer, the customer received two parcels. One from RedWood containing the Central Unit, the keyboard and the mouse. And an other one directly from Sony containing the MegaDisplay.\\
\par
SPECS HERE
\par
\pagebreak



\fullimage{next/NeXTcube_motherboard.png}
\par
\pagebreak



\drawing{NeXTCube_motherboard}{NextCube motherboard}
\pagebreak



\begin{figure}[H]
\centering
\scaledimage{1}{68030_blueprint.png}
\end{figure}
\pagebreak
\par


\begin{figure}[H]
\centering
\scaledimage{1}{68030_layout.png}
\end{figure}
\pagebreak
 


\begin{wrapfigure}[9]{r}{0.33\textwidth}
\centering
\scaledimage{0.33}{next/next-cube-system.png}
\end{wrapfigure}


Contrary to PCs which were a mess of wires, the NeXT Computer formed a chain. The mouse was connected to the keyboard, itself connected to the screen, connected to the Cube.\\


\par


If initially the NeXTComputer was acclaimed for its specs, there was a serious issue with the price. Market studies showed that students and researchers wanted a workstation priced at \$3,000. The NeXT Computer started at more than twice the ideal price at \$6,500. To make it worth, the optical drive which powered the bare configuration would turn out to be great for backup but way too slow for runtime. This made the "optional" \$3,500 SCSI hard-drive and absolute necessity which pushed the final price tag to \$10,000!\\
\par




\section{Line of Products}
With a factory tailed made to produce XXX units per week but week nearing XXX/month, NeXT upgraded its product line.\\
\par
In 1991 NeXT released four new products an upgraded NextComputer called NextCube. A smaller, flattened version of the NeXTCube which did not offer any expansion slots but with built-in color capability called NeXTStation. To offer color and video capture, the NeXTDimension expansion board.\\
\par
\newcolumntype{L}[1]{>{\hsize=#1\hsize\raggedright\arraybackslash}X}%
\newcolumntype{R}[1]{>{\hsize=#1\hsize\raggedleft\arraybackslash}X}%
 \begin{figure}[H]
\centering  
\begin{tabularx}{\textwidth}{ L{1.7}  R{0.3}  R{1.4}  R{0.8}  R{0.8}}
  \toprule
  \textbf{Name} &  \textbf{Year} & \textbf{Specs} & \textbf{\$\$\$} & \textbf{\$(2018)}   \\
  \toprule 
   NeXT Computer           & 1989 & 68030 25 Mhz & \$6,500 & \$12,938 \\
\toprule 
   NeXTStation             & 1991 & 68040 25 Mhz & \$4,995 & \$9,157 \\
   NeXTCube                & 1991 & 68040 25 Mhz & \$12 395 & \$21,171 \\
   NeXTDimension           & 1991 & i860  33 Mhz & \$3,995 & \$7,552 \\
   NeXTStation Color       & 1991 & 68040 25 Mhz & \$7,995 & \$14,656 \\
\toprule 
   NeXTCube Turbo          & 1992 & 68040 33 Mhz & \$10,000 & \$18,121 \\
   NextStation Turbo       & 1992 & 68040 33 Mhz & \$5995 & \$11,932 \\
   
   NeXTStation Turbo Color & 1992 & 68040 33 Mhz & \$8995 & \$17,904 \\
   \toprule
\end{tabularx}
\caption{Next products over the years\protect\footnotemark.}
\end{figure}
\par
\footnotetext{Source: kevra.org (Competing Hardware Comparisons) and https://simson.net/ref/NeXT/specifications.htm.}
\par
In 1992, they buffed up their entire line with Turbo version and what would become the Gold Standard at id Software: The NeXTStation Color Turbo.









\section{NeXTCube}

\begin{minipage}{\textwidth}
\scaledimage{0.5}{next/next-crt-top.png} \scaledimage{0.5}{next/next-cube-top.png}
\end{minipage}
\pagebreak


\section{NeXTStation}

\fullimage{cube_vs_slab.png}
\pagebreak

\par
\fullimage{next/nextstation.png}
\par
SPECS
\pagebreak





\section{NeXTDimension}
\fullimage{NeXTtv.png}
I think we used NextTV to grab the video from the camcorder.\\
I believe the video camera was a Hi8 Handycam Sony model.\\
\scaledimage{0.3}{SonyHi8.png}
\par

dual screen. Monster with no driver and no SDK.\\
\par 
\fq{I was sad that I never got dev tools for the i860 - I wanted to have it do the low level texture mapping primitives.}{John Carmack}

\pagebreak
\cfullimage{next/NextDimension_board.png}{NeXTDimension board}
\par
\drawing{NeXTDimension_motherboard}{}
\cfullimage{next/NeXTDimension.png}{NeXTDimension ad}
\par

\fq{There is one specification change. NeXT has eliminated the C-Cube Microsystems CL-550 JPEG chip from NeXTdimension. This is because our supplier, C-Cube Microsystems, has failed to deliver chips that meet their specifications.}{Felipe\_Fuster\at NeXT.COM}

\fq{The NeXTdimension ran a custom kernel which was designed to do soft realtime management of multiple threads within a single address space, provide demand paged virtual memory, and provide a source-compatible Mach API subset and full Mach messaging interface, along with a minimal UNIX system call API, just enough to implement the RenderMan back end and the PostScript device layer. The kernel was called "Graphics aCcelerator Kernel, or "GaCK". Yes, this was a jape at the funny capitalization of the company name. It was not Mach, or BSD, or Minix, or Linux.}{M Paquette}

\fq{Cray on a Chip\\
\par
The Intel 80860 was an impressive chip, able at top speed to perform close to 66 MFLOPS at 33 MHz in real applications, compared to a more typical 5 or 10 MFLOPS for other CPUs of the time. Much of this was marketing hype, and it never become popular, lagging behind most newer CPUs and Digital Signal Processors in performance.
The 860 has several modes, from regular scaler mode to a superscalar mode that executes two instructions per cycle and a user visible pipeline mode (instructions using the result register of a multi-cycle op would take the current value instead of stalling and waiting for the result). It can use the 8K data cache in a limited way as a small vector register (like those in supercomputers). The unusual cache uses virtual addresses, instead of physical, so the cache has to be flushed any time the page tables changes, even if the data is unchanged. Instruction and data busses are separate, with 4 G of memory, using segments. It also includes a Memory Management Unit for virtual storage.\\
\par
The 860 has thirty two 32 bit registers and thirty two 32 bit (or sixteen 64 bit) floating point registers. It was one of the first microprocessors to contains not only an FPU as well as an integer ALU, and also included a 3-D graphics unit (attached to the FPU) that supports lines drawing, Gouraud shading, Z-buffering for hidden line removal, and operations in conjunction with the FPU. It was also the first able to do an integer operation, and a (unique at the time) multiply and add floating point instruction, for the equivalent of three instructions, at the same time.\\
\par
However actually getting the chip at top speed usually requires using assembly language - using standard compilers gives it a speed closer to other processors. Because of this, it was used as a coprocessor, either for graphics, or floating point acceleration, like add in parallel units for workstations. Another problem with using the Intel 860 as a general purpose CPU is the difficulty handling interrupts. It is extensively pipelined, having as many as four pipes operating at once, and when an interrupt occurs, the pipes can spill and lose data unless complex code is used to clean up. Delays range from 62 cycles (best case) to 50 microseconds (almost 2000 cycles).}{John Bayko, University of Regina}
\par

\bu{Display PostScript and the NeXTdimension board}\\ 
\par
\fq{The Display PostScript system can be (very broadly) broken into two pieces, the PostScript interpreter and the device. The interpreter processes the language, and passes marking, imaging, and (in the NeXT version) compositing directions to the device layer.\\ 
\par
The device layer takes the high level marking, imaging, and compositing operations and (eventually) converts these to bitmap level operations. The Display PostScript system spends the majority of it's time down here.\\ 
\par
In the case of the NeXTdimension board, the device layer is implemented on the NeXTdimension board. Marking, imaging, and compositing operations are asynchronously transmitted to the NeXTdimension for processing while additional PostScript is interpreted on the 68K processor. A good degree of parallelism is achieved in normal operation. The NXPing() AppKit call is interpreted by the Display PostScript system as a request to synchronize the NeXTdimension, PostScript interpreter, and app, and will not return until PostScript rendering is complete on the NeXTdimension.\\}{M Paquette} 

\bu{Quick Renderman and the NeXTdimension board}\\
\par 
\fq{
Renderman uses a similar architecture. The 3DKit and RI C binding reside on the m68k (otherwise you couldn't link with them). Depending on the setup of the Renderman context, a RIB stream can be spooled to Photorealistic Renderman running on the host CPU (the m68K for black hardware), or to a Quick Renderman implementation loaded on demand into the Window Server. The Quick Renderman implementation in the Window Server may then, if the target window is on a NeXTdimension, forward the rendering operations to a Quick Renderman context running on the NeXTdimension board.}{M Paquette} 













\section{NeXTStep}
 Extract from brochure.

 Our collaboration with Higher Education provided the insight needed to visualize the seven breakthroughs that would ultimately define the NeXT Computer:\\
\begin{enumerate}
\item A new architecture optimized for total system throughput, not just individual component benchmarks.

\item A pioneering technology for vast and reliable storage, opening the door for new ways to access and use information.

\item Built-in CD-quality sound, allowing sound to be integrated into applications that are used every day.

\item A unified imaging system - Display PostScript - for both the display and the printer. So what you see on the screen is unequivocally what you get on paper.

\item An intuitive interface that gives everyone access to UNIX, with all of its power for networking and multitasking.

\item A multimedia mail system that enables communication combining text, graphics, and voice.

\item A new development environment that dramatically cuts the time it takes to create and customize software.
\end{enumerate}

\fullimage{next/nextstep.png}
\par
Trivia: Macosx window capture icon is a remanend of NextStep.\\
\par
\footnote{Source: "NeXT vs Sun: A world of a difference", 1991 promotional video}

















\section{NeXT at id Software}
While some hated it.\\
\par
\fq{Develop for it? I'll piss on it.}{- Bill Gates (when InfoWorld asked if Microsoft would develop for the NeXT computer).}\\
\par
Some loved it.\\
\par
\fq{
When id Software was stationed in Madison, Wisconsin during the winter of 1991, most of us were gone for the Christmas holiday - except John Carmack. John's present, which he bought with \$11,000 of his own money, procured by walking through the snow and ice to remove from the bank, arrived during the holiday and he spent the whole time learning as much as he could about the computer and started working on vector quantization algorithms for compressing graphics. His test graphic was a 256-color screen from King's Quest 5.\\
\par
 After his research was done it was agreed that the entire company needed to develop our next game on NeXTSTEP.}{John Romero, rome.ro, December 20, 2006}\\
The NeXTStation had impressed id Software so much they decided to make it the foundation of the development of DOOM, relegating PCs to running the game. Only Deluxe Paint would be used on PC, eveything else, from programmation to level production, asset packaging and book writing would be done on NeXT.


\par
In 1990, NeXT unveiled faster workstations running at 25MHz on the Motorola 68040. Affectionately called "the slab," the NeXTstation had a black and white display while the NeXTstation Color displayed 4,096 colors from a palette of 16 million colors. A new version of the Cube offered a 32-bit true-color display.
Spring of 1992 brought about the end of NeXT's optical disk, but NeXT also introduced upgraded workstations, the NeXTstation Turbo and Turbo Color running at 33MHz. NeXT also began offering a color printer and a standard CD-ROM drive. Foreshadowing its exit from the hardware business, the company announced it had begun working on NeXTstep for Intel.\\



\par


NeXT was also used for Doom Battle Book and Wolfestein3D hint book.\\
NeXT's system software was designed to rival the best offerings of the Macintosh and PC: NeXT used the rock-solid UNIX operating system and added its own elegant, proprietary graphical user interface.  NeXT was also the first computer company to ship a built-in 256 MB magneto-optical storage medium.  Boasting a high-resolution display, built-in Ethernet, CD-quality sound, and multimedia e-mail, the NeXT Computer was packaged in a stunning one-foot by one-foot black magnesium cube.\\
\par


\par
\fq{
id's first NeXT hardware was all black - both Cubes and Stations. We upgraded through the years to the Turbo model then to other hardware like the HP Gecko and then Intel hardware at the end. We were building fat binaries of the tools for all 3 processors in the office - one .app file that had code for all 3 processors in it and executed the right code depending on which machine you ran it on. All our data was stored on a Novell 3.11 server and we constantly used the NeXTSTEP Novell gateway object to transparently copy our files to and from the server as if it was a local NTFS drive. This was back in 1993!\\
\par
In fact, with the superpower of NeXTSTEP, one of the earliest incarnations of DoomEd had Carmack in his office, me in my office, DoomEd running on both our computers and both of us editing one map together at the same time. I could see John moving entities around on my screen as I drew new walls. Shared memory spaces and distributed objects. Pure magic.\\
\par
We wrote all of DOOM and Quake's code on NeXTSTEP. We debugged the code in NeXTSTEP with DOOM and Quake's 320x200 VGA screen drawing in a little Interceptor window while the rest of the screen was used for debugging code. When all the code ran without bugs we cross-compiled it for the Intel processor on NeXTSTEP then turned over to our Intel DOS computers, copied the EXE and just ran the game. The DOS4GW DOS-Extender loaded up and the game ran. It was that easy.}{}

\fq{
One funny and strange note: I'm left-handed but while using NeXTSTEP I used a right-hand mouse. Then when I ran the game on my DOS computer I switched hands and played left-handed. Windows is not yet hardcore enough for me to switch to right-handed mice.\\
\par
I'll bet you didn't know that DOOM, DOOM II and Quake weren't the only games developed on NeXTSTEP. When I got Raven Software to agree to develop Heretic for us I had them buy several Epson NeXT computers (Intel based) and I flew up to Madison, WI to get them all set up and teach them how to develop the game with our tools and engine. It was a great time I'll never forget - seeing their team get excited about the power of the new environment and that they got the game developed and released in under a year. They signed on for another title and developed Hexen on NeXTSTEP as well.\\
\par
Back at home base we were busy with the beginnings of Quake - developed on NeXTSTEP of course but we also had another game that we were having developed using our tools and tech by Rogue Entertainment - the super fun, one of a kind, action RPG that was Strife. Unfortunately most people hadn't heard of the game because the publisher we sold the game to went out of business upon its launch. If you can find this game you should really play it - it's quite a ride.\\
\par
As I was leaving id Software in August 1996 the move to the Windows 32 platform was underway. John Carmack was porting our QuakeEd editor over to Win32 and preparing for a NeXT-less future. Several short months later NeXT made their fateful move over to Apple and a new era was begun as Steve Jobs set about changing the future. Again.\\
\par
Up to that point I had spent 15 years of my life working on computers that Steve Jobs was involved in bringing to the world. First the Apple II+, then the IIe, the IIgs and finally NeXT. Maybe someday I'll get one of those kickass iMacs.}{John Romero, rome.ro, December 20, 2006}
\par
Retrospect by john carmack, answering the question "Why was Doom developed on a NeXT?":\\
\par

\fq{No regrets at all!\\
\par
I bought our first NeXT (a ColorStation) just out of personal interest. Jason Blochowiak had talked to me about the advantages of Unix based systems from his time at college, and I was interested in seeing what Steve Job’s next big thing was. It is funny to look back – I can remember honestly wondering what the advantages of a real multi process development environment would be over the DOS and older Apple environments we were using. Actually using the NeXT was an eye opener, and it was quickly clear to me that it had a lot of tangible advantages for us, so we moved everything but pixel art (which was still done in Deluxe Paint on DOS) over. Using Interface Builder for our game editors was a NeXT unique advantage, but most Unix systems would have provided similar general purpose software development advantages (the debugger wasn’t nearly as good as Turbo Debugger 386, though!). Kevin Cloud even did our game manuals, starting with Wolfenstein 3D, in Framemaker on a NeXT.\\
\par
This was all in the context of DOS or Windows 3.x; it was revelatory to have a computer system that didn’t crash all the time. By the time Quake 2 came around, Windows NT was in a similar didn’t-crash-all-the-time state, it had hardware accelerated OpenGL, and Visual Studio was getting really good, so I didn’t feel too bad moving over to it. At that transition point I did evaluate most of the other Unix workstations, and didn’t find a strong enough reason not to go with Microsoft for our desktop systems.\\
\par
Over the entire course of Doom and Quake 1’s development we probably spent \$100,000 on NeXT computers, which isn’t much at all in the larger scheme of development. We later spent more than that on Unix SMP server systems (first a quad Alpha, then an eventually 16-way SGI system) to run the time consuming lighting and visibility calculations for the Quake series. I remember one year looking at the Top 500 supercomputer list and thinking that if we had expanded our SGI to 32 processors, we would have just snuck in at the bottom.}{John Carmack, quora.com}
\par
\par
\fq{The NeXTCube was purchased in December 1991 and was the only NeXT Computer we had until December 1992 when we decided we would develop DOOM with them so we bought 3 NeXTStations: mine, John Carmack’s new one, Tom Hall’s. John C’s original NeXTCube was the computer used to scan the clay models, gun toys, and latex models.}{John Romero}\\
\par
comp.sys.next.advocacy ›
DOOM: NeXTstep's Most Successful App


\section{After NeXT}
NeXT stopped manufacturing hardware in 1993 to become a software-only vendor, selling NeXTSTEP as a combination operating system and object-oriented development environment. NeXTstep for Intel became a popular product among large companies and especially financial institutions for rapidly developing and deploying custom software. \\
\par
Apple Computer bought NeXT in 1996 after its own efforts to upgrade the Macintosh operating system failed.  After the sale, Steve Jobs first began working as an advisor but was later appointed acting-CEO, and then finally CEO of the company.\\
\par
NeXTSTEP lives on as the heart of Mac OS X. It is likely some of the code that once ran on NeXTStep and contributed to \doom now runs on the many mac computer across the world.\\


